\documentclass[11pt]{article}

\usepackage[a4paper,margin=1in]{geometry}
\usepackage{amsmath,amssymb,amsthm}
\usepackage{bm}
\usepackage{mathtools}
\usepackage{enumitem}
\usepackage{booktabs}
\usepackage{xcolor}
\usepackage{microtype}
\usepackage[colorlinks=true,linkcolor=NavyBlue,urlcolor=NavyBlue,citecolor=NavyBlue]{hyperref}
\usepackage[dvipsnames]{xcolor}
\usepackage{titlesec}
\usepackage{fancyhdr}
\usepackage{tcolorbox}
\tcbuselibrary{skins,breakable}

% --- light styling ---------------------------------------------------
\titleformat{\section}{\Large\bfseries\color{MidnightBlue}}{\thesection}{0.6em}{}
\titleformat{\subsection}{\large\bfseries\color{MidnightBlue!85!black}}{\thesubsection}{0.6em}{}
\pagestyle{fancy}\fancyhf{}
\lhead{\small FRACMATH --- Theory Manual}
\rhead{\small CDM, modified von Mises, Newton--Raphson}
\cfoot{\thepage}
\renewcommand{\headrulewidth}{0.3pt}

\newtcolorbox{algobox}[1][]{
  enhanced, breakable, colback=Gray!6, colframe=MidnightBlue!70,
  boxrule=0.6pt, arc=2pt, left=8pt, right=8pt, top=6pt, bottom=6pt, #1}

% --- shorthand -------------------------------------------------------
\newcommand{\eps}{\bm{\varepsilon}}
\newcommand{\sig}{\bm{\sigma}}
\newcommand{\C}{\mathbb{C}}
\newcommand{\Id}{\bm{I}}
\newcommand{\dev}{\mathrm{dev}}
\newcommand{\tr}{\mathrm{tr}}
\newcommand{\eqs}{\tilde{\varepsilon}}
\newcommand{\dd}{\,\mathrm{d}}
\newcommand{\norm}[1]{\left\lVert#1\right\rVert}

\title{\textbf{Damage Mechanics in FRACMATH}\\[2pt]
\large A worked derivation of the scalar CDM model,\\
the modified von Mises equivalent strain, and the Newton--Raphson solver}
\author{Jaykumar Mavani \quad\textbar\quad Madura Pathirage\\
\small Dept.\ of Civil, Construction \& Environmental Engineering, University of New Mexico}
\date{\today}

\begin{document}
\maketitle
\thispagestyle{fancy}

\begin{abstract}
\noindent
This manual derives the constitutive model behind FRACMATH, step by step, with nothing
swept under the rug. It covers the scalar isotropic damage law, the modified von Mises
equivalent strain of de Vree and co-workers, exponential softening, crack band
regularization through Oliver's characteristic length, and the modified Newton--Raphson
scheme the code actually solves with. Where a formula looks like it was pulled from thin
air, we check it against a uniaxial test so you can see why it has the shape it does.
The aim is that a reader can follow every line from the free energy down to the element
tangent and reproduce the model in any FE code, not only this one.
\end{abstract}

\tableofcontents
\bigskip

% =====================================================================
\section{Scope and notation}
% =====================================================================

We model fracture in concrete and rock with one scalar damage variable. That is a
deliberate choice. A single $\omega$ cannot represent crack closure or direction-dependent
stiffness loss, but for monotonic loading it is accurate, cheap, and easy to debug, and it
is the workhorse of structural-scale fracture simulation. The whole formulation lives in
small strains and rate independence.

Tensors are bold. Second-order tensors use a single bold letter ($\eps$, $\sig$), the
fourth-order elastic tensor is $\C_0$, and $\Id$ is the second-order identity. A colon is a
double contraction, $\C_0:\eps$ has components $C_{ijkl}\varepsilon_{kl}$, and a dyadic
$\bm a \otimes \bm b$ has components $a_{ij}b_{kl}$. The strain invariants we need are
%
\begin{equation}
I_1 = \tr\eps = \varepsilon_{kk},
\qquad
\bm e = \eps - \tfrac{1}{3} I_1 \Id,
\qquad
J_2 = \tfrac{1}{2}\,\bm e : \bm e ,
\label{eq:invariants}
\end{equation}
%
where $\bm e$ is the deviatoric strain. The factor of one half in $J_2$ is not cosmetic; it
is exactly what makes the equivalent strain reduce to the axial strain under uniaxial
tension, and getting it wrong (a stray $1/6$, say) shifts the whole damage threshold. We
verify that in Section~\ref{sec:eqstrain}.

% =====================================================================
\section{Damage kinematics and a short thermodynamic argument}
% =====================================================================

\subsection{Effective stress and the stiffness reduction}

Start from the physical picture. A representative volume of material carries microcracks
and voids. Let $\omega \in [0,1]$ be the fraction of a cross-section that has been lost to
those defects, so $\omega=0$ is virgin material and $\omega=1$ is fully broken. The load
that used to pass through the whole section now passes through the remaining ligament, so
the \emph{effective} stress on the intact material is
%
\begin{equation}
\bar{\sig} = \frac{\sig}{1-\omega}.
\end{equation}
%
The intact material is still linear elastic, $\bar{\sig} = \C_0 : \eps$, which gives the
nominal stress carried by the damaged solid:
%
\begin{equation}
\boxed{\;\sig = (1-\omega)\,\C_0 : \eps\;}
\label{eq:stress}
\end{equation}
%
This is the entire stress--strain law. Everything that follows is about one thing: how
$\omega$ grows.

\subsection{Why the law is allowed (the dissipation check)}

It is worth confirming that \eqref{eq:stress} does not secretly create energy. Take the
Helmholtz free energy per unit volume as the stored elastic energy scaled by the surviving
stiffness,
%
\begin{equation}
\psi(\eps,\omega) = (1-\omega)\,\tfrac{1}{2}\,\eps : \C_0 : \eps = (1-\omega)\,\psi_0(\eps).
\end{equation}
%
The Clausius--Duhem inequality for an isothermal process reads
$\mathcal{D} = \sig:\dot{\eps} - \dot\psi \ge 0$. Differentiating $\psi$,
%
\begin{equation}
\dot\psi = (1-\omega)\,\C_0:\eps : \dot\eps \;-\; \psi_0\,\dot\omega
        = \sig : \dot\eps \;-\; \psi_0\,\dot\omega,
\end{equation}
%
so the dissipation collapses to
%
\begin{equation}
\mathcal{D} = \psi_0(\eps)\,\dot\omega \;\ge\; 0.
\end{equation}
%
Since $\psi_0 \ge 0$ always, the model is thermodynamically admissible as long as damage
never decreases, $\dot\omega \ge 0$. That single constraint, no healing, is what the
history variable in the next section enforces.

% =====================================================================
\section{The modified von Mises equivalent strain}
\label{sec:eqstrain}
% =====================================================================

Damage should grow when the strain state is severe enough, but ``severe'' has to be a
single number before we can compare it to a threshold. That number is the equivalent
strain $\eqs$. Concrete is far weaker in tension than in compression, so $\eqs$ has to be
asymmetric: a given tensile strain must count for much more than the same strain in
compression. The modified von Mises measure of de Vree, Brekelmans and van Gils builds that
asymmetry in through the strength ratio
%
\begin{equation}
k = \frac{f_c}{f_t}, \qquad k \approx 10 \text{ for concrete}.
\end{equation}
%
The measure is
%
\begin{equation}
\eqs
= \underbrace{\frac{k-1}{2k(1-2\nu)}\,I_1}_{\text{volumetric, signed}}
+ \frac{1}{2k}\sqrt{\Bigl(\frac{k-1}{1-2\nu}\Bigr)^{2} I_1^{2}
        + \frac{12k}{(1+\nu)^{2}}\,J_2}.
\label{eq:eqstrain}
\end{equation}
%
The first term carries the sign of the volumetric strain, so tension and compression are
treated differently from the start. The term under the root is positive either way and
supplies the deviatoric (shape-change) contribution. The grouping of constants looks
arbitrary until you push two load cases through it.

\subsection{Why these exact constants: the uniaxial tension check}

Put the material in uniaxial tension with axial strain $\varepsilon>0$. The lateral
strains are $-\nu\varepsilon$, so $\eps = \mathrm{diag}(\varepsilon,-\nu\varepsilon,-\nu\varepsilon)$.
The first invariant is
%
\begin{equation}
I_1 = \varepsilon(1-2\nu).
\end{equation}
%
The deviatoric components come out to
%
\begin{equation}
e_{11} = \tfrac{2}{3}(1+\nu)\varepsilon, \qquad
e_{22}=e_{33} = -\tfrac{1}{3}(1+\nu)\varepsilon,
\end{equation}
%
and the deviatoric invariant is
%
\begin{equation}
J_2 = \tfrac{1}{2}\bigl(e_{11}^2+e_{22}^2+e_{33}^2\bigr)
    = \tfrac{1}{2}\cdot\tfrac{6}{9}(1+\nu)^2\varepsilon^2
    = \tfrac{1}{3}(1+\nu)^2\varepsilon^2.
\end{equation}
%
Now feed these into the two pieces of \eqref{eq:eqstrain}. The deviatoric term inside the
root cancels the $(1+\nu)^2$ cleanly,
%
\begin{equation}
\frac{12k}{(1+\nu)^2}J_2 = \frac{12k}{(1+\nu)^2}\cdot\tfrac{1}{3}(1+\nu)^2\varepsilon^2 = 4k\,\varepsilon^2,
\end{equation}
%
and the volumetric term inside the root cancels the $(1-2\nu)$,
%
\begin{equation}
\Bigl(\frac{k-1}{1-2\nu}\Bigr)^2 I_1^2
= \Bigl(\frac{k-1}{1-2\nu}\Bigr)^2 \varepsilon^2(1-2\nu)^2 = (k-1)^2\varepsilon^2.
\end{equation}
%
So the root simplifies to a perfect square:
%
\begin{equation}
\sqrt{(k-1)^2\varepsilon^2 + 4k\varepsilon^2}
= \varepsilon\sqrt{k^2+2k+1} = \varepsilon(k+1).
\end{equation}
%
The leading term is $\frac{k-1}{2k(1-2\nu)}\varepsilon(1-2\nu) = \frac{k-1}{2k}\varepsilon$, and putting
both together,
%
\begin{equation}
\eqs = \frac{k-1}{2k}\varepsilon + \frac{1}{2k}\varepsilon(k+1)
     = \frac{\varepsilon}{2k}\bigl[(k-1)+(k+1)\bigr] = \varepsilon.
\end{equation}
%
That is the point of the constants. Under uniaxial tension the equivalent strain is just
the axial strain, so the threshold $\kappa_0$ is read straight off a tension test with no
correction factor.

\subsection{And the compression check}

Repeat with uniaxial compression, axial strain $-\varepsilon$ ($\varepsilon>0$). Now
$I_1=-\varepsilon(1-2\nu)$, the first term flips sign, the root is unchanged, and
%
\begin{equation}
\eqs = -\frac{k-1}{2k}\varepsilon + \frac{k+1}{2k}\varepsilon
     = \frac{\varepsilon}{2k}\bigl[-(k-1)+(k+1)\bigr] = \frac{\varepsilon}{k}.
\end{equation}
%
Compression reaches the same $\eqs$ only at $k$ times the strain. Damage therefore starts
in compression at roughly $k$ times the tensile load, which is exactly the strength ratio
$f_c/f_t$ we asked for. The measure does what concrete does.

% =====================================================================
\section{Damage growth: loading function and history}
% =====================================================================

Damage tracks the worst strain the point has ever felt. Define a history variable $\kappa$
and a loading function
%
\begin{equation}
f(\eqs,\kappa) = \eqs - \kappa .
\end{equation}
%
The state must satisfy the Kuhn--Tucker (loading--unloading) conditions
%
\begin{equation}
f \le 0, \qquad \dot\kappa \ge 0, \qquad \dot\kappa\, f = 0 .
\end{equation}
%
Read them in plain terms. The strain state can never sit outside the current envelope
($f\le 0$); the memory only ever grows ($\dot\kappa\ge0$); and it grows only while the
state is pushing on the envelope ($\dot\kappa f=0$). Integrated over a monotonic-in-memory
history this is simply
%
\begin{equation}
\kappa(t) = \max\Bigl(\kappa_0,\ \max_{\tau\le t}\eqs(\tau)\Bigr),
\label{eq:history}
\end{equation}
%
with the initial threshold
%
\begin{equation}
\kappa_0 = \frac{f_t}{E}.
\end{equation}
%
Below $\kappa_0$ the material is elastic and $\omega=0$. The instant $\eqs$ first exceeds
$\kappa_0$, damage starts and never reverses, which is the no-healing condition the
dissipation argument demanded.

% =====================================================================
\section{Exponential softening}
% =====================================================================

We need a rule $\omega(\kappa)$ that climbs from $0$ at $\kappa_0$ toward $1$ as $\kappa$
grows. FRACMATH uses the exponential law
%
\begin{equation}
\omega(\kappa) = 1 - \frac{\kappa_0}{\kappa}\,
\exp\!\left[-\frac{\kappa-\kappa_0}{\varepsilon_f-\kappa_0}\right],
\qquad \kappa \ge \kappa_0 .
\label{eq:softening}
\end{equation}
%
Two limits check out by inspection: at $\kappa=\kappa_0$ the exponential is $1$ and the
prefactor is $1$, so $\omega=0$; as $\kappa\to\infty$ both the $\kappa_0/\kappa$ factor and
the exponential go to zero, so $\omega\to1$. The parameter $\varepsilon_f$ sets how fast the
material lets go, and we will pin it down from fracture energy in the next section.

A detail that matters for the solver: the uniaxial stress during monotonic loading takes a
clean closed form. With $\varepsilon=\kappa$ during loading and $E\kappa_0=f_t$,
%
\begin{equation}
\sigma = (1-\omega)E\kappa
       = E\kappa_0\exp\!\left[-\frac{\kappa-\kappa_0}{\varepsilon_f-\kappa_0}\right]
       = f_t\,\exp\!\left[-\frac{\kappa-\kappa_0}{\varepsilon_f-\kappa_0}\right].
\label{eq:softstress}
\end{equation}
%
The stress decays exponentially from the tensile strength. We use this in a moment to get
the dissipated energy in one integral.

We will also need the slope of the damage law for the tangent. Differentiating
\eqref{eq:softening} with the product rule,
%
\begin{equation}
\frac{\dd\omega}{\dd\kappa}
= \frac{\kappa_0}{\kappa}\exp\!\left[-\frac{\kappa-\kappa_0}{\varepsilon_f-\kappa_0}\right]
\left(\frac{1}{\kappa} + \frac{1}{\varepsilon_f-\kappa_0}\right),
\qquad \kappa \ge \kappa_0 .
\label{eq:domega}
\end{equation}
%
It is positive throughout, as a softening law must be, since every factor is positive.

% =====================================================================
\section{Crack band regularization and the calibration of $\varepsilon_f$}
% =====================================================================

Plain softening has a famous problem. Refine the mesh and the damage localizes into one
element, the dissipated energy shrinks with the element size, and the answer never
converges. The crack band method fixes this by making the local softening branch depend on
the element size so that the energy released per unit crack \emph{area} stays fixed at the
material fracture energy $G_F$.

\subsection{Energy under the softening curve}

The energy dissipated per unit volume in a fully softened material point is the area under
the uniaxial stress--strain curve. Split it into the elastic triangle up to $\kappa_0$ and
the softening tail from \eqref{eq:softstress}:
%
\begin{align}
g_f &= \underbrace{\tfrac{1}{2} f_t \kappa_0}_{\text{elastic}}
     + \underbrace{\int_{\kappa_0}^{\infty} f_t
       \exp\!\left[-\frac{\kappa-\kappa_0}{\varepsilon_f-\kappa_0}\right]\dd\kappa}_{\text{softening tail}}
     \\[2pt]
    &= \tfrac{1}{2} f_t \kappa_0 + f_t(\varepsilon_f-\kappa_0)
     = f_t\!\left(\varepsilon_f - \tfrac{1}{2}\kappa_0\right).
\label{eq:gf}
\end{align}
%
The tail integral is elementary: the integral of a decaying exponential is just its decay
length $(\varepsilon_f-\kappa_0)$ times the starting value $f_t$.

\subsection{Tie the volume energy to the area energy}

Crack band regularization says the energy per unit \emph{area}, $G_F$, equals the energy
per unit \emph{volume}, $g_f$, smeared over a band one characteristic length $h$ wide:
%
\begin{equation}
G_F = h\,g_f .
\end{equation}
%
Substitute \eqref{eq:gf} and solve for the only free parameter:
%
\begin{equation}
\boxed{\;\varepsilon_f = \frac{G_F}{h\,f_t} + \frac{\kappa_0}{2}\;}
\label{eq:epsf}
\end{equation}
%
Each element gets its own $\varepsilon_f$ from its own $h$. A big element softens faster
(smaller $\varepsilon_f$), a small element softens slower, and the product $h\,g_f$ comes
out to $G_F$ either way. That is what makes the peak load and the total dissipation
independent of the mesh.

\subsection{The characteristic length: Oliver's projection}

The remaining question is what $h$ to use. For a linear triangle, FRACMATH uses Oliver's
projection,
%
\begin{equation}
h = \frac{2A_e}{\ell_{\max}},
\end{equation}
%
where $A_e$ is the element area and $\ell_{\max}$ its longest edge. Geometrically this is
the width of the element measured across the band, i.e.\ the height of the triangle on its
longest side. The longest edge is taken as the proxy for the crack direction, so $h$ is the
distance the crack has to traverse. For 3D elements the same idea uses the element volume
and the projected area of the dominant face. The practical payoff, and the reason the
torsion benchmark works at all, is that $h$ follows the local crack orientation instead of
being locked to a fixed mesh diagonal.

% =====================================================================
\section{From the continuum to a discrete residual}
% =====================================================================

The equilibrium statement is the principle of virtual work: for every admissible virtual
displacement $\delta\bm u$,
%
\begin{equation}
\int_\Omega \delta\eps : \sig \dd V
= \int_\Omega \delta\bm u\cdot\bm b \dd V
+ \int_{\Gamma_t}\delta\bm u\cdot\bm t \dd S .
\end{equation}
%
Interpolate the displacement with shape functions, $\bm u = \bm N\bm d$, so that strain is
$\eps = \bm B\bm d$ with $\bm B$ the usual strain--displacement matrix. The virtual work
statement turns into a set of nonlinear algebraic equations in the nodal displacements
$\bm d$:
%
\begin{equation}
\bm r(\bm d) = \bm f^{\text{int}}(\bm d) - \bm f^{\text{ext}} = \bm 0,
\qquad
\bm f^{\text{int}} = \int_\Omega \bm B^{\!\top}\sig \dd V .
\label{eq:residual}
\end{equation}
%
The internal force is nonlinear in $\bm d$ because $\sig$ depends on $\omega$, which depends
on the strain through $\eqs$ and the history. We solve \eqref{eq:residual} under
displacement control, prescribing the loaded degrees of freedom and iterating on the free
ones.

% =====================================================================
\section{The Newton--Raphson solve}
% =====================================================================

\subsection{The consistent tangent, derived in full}

Newton's method linearizes the residual: given an iterate $\bm d^{(i)}$ with residual
$\bm r^{(i)}$, solve $\bm K_T\,\delta\bm d = -\bm r^{(i)}$ and update
$\bm d^{(i+1)} = \bm d^{(i)} + \delta\bm d$. The tangent is
$\bm K_T = \partial \bm f^{\text{int}}/\partial\bm d = \int_\Omega \bm B^{\!\top}
\bm D_T \bm B \dd V$, so the whole problem reduces to the material tangent
$\bm D_T = \dd\sig/\dd\eps$.

Differentiate the stress law \eqref{eq:stress}. During active loading the history follows
the equivalent strain, $\kappa=\eqs$ and $\dd\kappa=\dd\eqs$, so by the product rule
%
\begin{equation}
\dd\sig
= (1-\omega)\,\C_0:\dd\eps
\;-\; \frac{\dd\omega}{\dd\kappa}\,(\C_0:\eps)\,\dd\eqs .
\end{equation}
%
The first term is the secant stiffness acting on the strain increment. The second is the
loss of stiffness as damage advances. Writing $\dd\eqs = \dfrac{\partial\eqs}{\partial\eps}:\dd\eps$,
the consistent tangent is
%
\begin{equation}
\boxed{\;
\bm D_T = (1-\omega)\,\C_0
\;-\; \frac{\dd\omega}{\dd\kappa}\,
\bigl(\C_0:\eps\bigr)\otimes\frac{\partial\eqs}{\partial\eps}
\;}
\label{eq:tangent}
\end{equation}
%
The first part is symmetric. The second is a rank-one dyad and is generally not symmetric,
which is the source of the unsymmetric, sometimes ill-conditioned tangents that make full
Newton stall on the steep part of the softening branch.

\subsection{The gradient of the equivalent strain}

The dyad in \eqref{eq:tangent} needs $\partial\eqs/\partial\eps$. Write \eqref{eq:eqstrain}
compactly with
%
\begin{equation}
a = \frac{k-1}{2k(1-2\nu)}, \quad
b = \frac{k-1}{1-2\nu}, \quad
c = \frac{12k}{(1+\nu)^2}, \quad
S = \sqrt{b^2 I_1^2 + c\,J_2},
\end{equation}
%
so that $\eqs = a I_1 + \tfrac{1}{2k}S$. Using $\partial I_1/\partial\eps = \Id$ and
$\partial J_2/\partial\eps = \bm e$ (which follows from $J_2=\tfrac12\bm e:\bm e$ and the
deviatoric projection),
%
\begin{equation}
\frac{\partial\eqs}{\partial\eps}
= a\,\Id
+ \frac{1}{2k}\cdot\frac{1}{2S}\Bigl(2b^2 I_1\,\Id + c\,\bm e\Bigr)
= a\,\Id + \frac{b^2 I_1}{2kS}\,\Id + \frac{c}{4kS}\,\bm e .
\label{eq:deqs}
\end{equation}
%
As $S\to 0$ near a stress-free state this expression is singular, which is harmless because
damage is inactive there and the dyad term is never assembled; the code guards the divide
with the loading check anyway.

\subsection{What FRACMATH actually solves: the modified (secant) scheme}

Full Newton with \eqref{eq:tangent} converges quadratically but is fragile once the dyad
term turns the tangent unsymmetric and nearly singular. FRACMATH instead uses the
\emph{modified} Newton--Raphson scheme: keep only the symmetric secant part,
%
\begin{equation}
\bm D_T \approx (1-\omega)\,\C_0,
\qquad
\bm K_T = \int_\Omega \bm B^{\!\top}(1-\omega)\,\C_0\,\bm B \dd V .
\label{eq:secant}
\end{equation}
%
This drops the rank-one term entirely. The trade is deliberate: convergence slows from
quadratic to linear, but the tangent stays symmetric positive definite, the factorization
is stable, and it can be reused across iterations within a step. The residual
\eqref{eq:residual} still uses the true damaged stress, so the converged solution is the
exact one; only the path to it is slower. On the benchmarks this buys robustness through
the snap-back-prone softening regime at a cost of a few extra iterations per step, which
the displacement-controlled stepping absorbs.

\subsection{Solver loop}

\begin{algobox}
\textbf{Modified Newton--Raphson, one displacement-controlled step}
\begin{enumerate}[leftmargin=1.4em,itemsep=2pt]
\item Apply the prescribed displacement increment to the loaded DOFs; freeze the supports.
\item \textbf{For} iteration $i=0,1,2,\dots$ until $\norm{\bm r}\le \texttt{tol}$:
  \begin{enumerate}[leftmargin=1.4em,itemsep=1pt]
  \item For every element, compute $\eps=\bm B\bm d$, then $I_1$, $J_2$, and
        $\eqs$ from \eqref{eq:eqstrain}.
  \item Update history $\kappa=\max(\kappa,\eqs)$ from \eqref{eq:history};
        if $\kappa>\kappa_0$, update $\omega$ from \eqref{eq:softening} with the
        element's $\varepsilon_f$ from \eqref{eq:epsf}.
  \item Assemble $\sig=(1-\omega)\C_0:\eps$ and the internal force
        $\bm f^{\text{int}}=\int\bm B^{\!\top}\sig\dd V$.
  \item Form the residual $\bm r = \bm f^{\text{int}}-\bm f^{\text{ext}}$ on the free DOFs.
  \item Assemble the secant tangent $\bm K_T$ from \eqref{eq:secant};
        solve $\bm K_T\,\delta\bm d = -\bm r$ with sparse \texttt{backslash} (UMFPACK).
  \item Update $\bm d \leftarrow \bm d + \delta\bm d$.
  \end{enumerate}
\item Commit the converged $\kappa$ and $\omega$ as the history for the next step.
\end{enumerate}
\end{algobox}

All of step~2(a)--(c) runs over every element in one batched tensor operation
(\texttt{pagemtimes} in MATLAB), which is where the vectorized speedup comes from. The
sparse solve in 2(e) dominates the wall-clock time, as the timing table in the JOSS paper
shows.

% =====================================================================
\section{Symbol reference}
% =====================================================================

\begin{center}
\renewcommand{\arraystretch}{1.25}
\begin{tabular}{@{}cl@{\hskip 2em}cl@{}}
\toprule
Symbol & Meaning & Symbol & Meaning \\
\midrule
$\omega$        & scalar damage, $[0,1]$       & $E,\nu$         & Young's modulus, Poisson ratio \\
$\sig,\eps$     & stress, strain               & $\C_0$          & virgin elastic stiffness \\
$\eqs$          & equivalent strain            & $\kappa$        & history variable \\
$I_1$           & first strain invariant       & $J_2$           & second deviatoric invariant \\
$\bm e$         & deviatoric strain            & $k=f_c/f_t$     & strength ratio \\
$f_t,f_c$       & tensile, compressive strength& $\kappa_0=f_t/E$& damage onset threshold \\
$\varepsilon_f$ & softening parameter          & $G_F$           & fracture energy \\
$g_f$           & energy per unit volume       & $h$             & characteristic length \\
$A_e,\ell_{\max}$& element area, longest edge  & $\bm B$         & strain--displacement matrix \\
$\bm r$         & residual                     & $\bm K_T$       & tangent stiffness \\
$\bm D_T$       & material tangent             & $\bm d$         & nodal displacements \\
\bottomrule
\end{tabular}
\end{center}

\bigskip
\noindent\textbf{Primary references.} de Vree, Brekelmans \& van Gils (1995) for the
equivalent strain; Ba\v{z}ant \& Oh (1983) for crack band theory; Oliver (1989) for the
characteristic-length projection; Peerlings et al.\ (1996) and Jir\'asek (2004) for the
damage formulation and its numerics. Full citations are in the repository
\texttt{paper.bib}.

\end{document}
